% =========================
% report/evaluation.tex
% =========================

\subsection{هدف ارزیابی}
ارزیابی باید سه بُعد را پوشش دهد:
\begin{enumerate}
  \item \textbf{درستی هندسی:} آیا $\mathbf{H}$ واقعاً الگو را روی عکس درست قرار می‌دهد؟
  \item \textbf{پایداری:} عملکرد تحت تغییر نور/زاویه/شلوغی چقدر افت می‌کند؟
  \item \textbf{کارایی:} زمان اجرا برای کاربرد عملی مناسب است یا خیر؟
\end{enumerate}

\subsection{معیارهای کمی اصلی}
\paragraph{خطای بازفرافکنی (پیکسل).}
اگر نقاط زمین‌حقیقت $(\mathbf{p}_k^T,\mathbf{p}_k^Q)$ موجود باشد:
\[
e_k = \left\|\mathbf{p}_k^Q - \pi(\mathbf{H}\tilde{\mathbf{p}}_k^T)\right\|_2.
\]
اینجا $\|\cdot\|_2$ همان نرم اقلیدسی است:
\[
\|[a,b]^T\|_2=\sqrt{a^2+b^2}.
\]
\textbf{تفسیر:} $e_k$ یعنی نقطهٔ پیش‌بینی‌شده توسط هموگرافی چند پیکسل با نقطهٔ واقعی فاصله دارد.

\paragraph{میانه/میانگین خطا در هر تصویر.}
برای هر تصویر، از مجموعه خطاها (مثلاً ۱۰ نقطه) یک خلاصه مثل \textbf{میانه} یا \textbf{میانگین} گرفته می‌شود. میانه نسبت به نقاط پرت مقاوم‌تر است.

\paragraph{نسبت درون‌نهشتی‌ها.}
اگر $N_m$ تعداد تطبیق‌ها و $N_i$ تعداد درون‌نهشتی‌ها باشد:
\[
\text{InlierRatio} = \frac{N_i}{N_m}.
\]
\textbf{تفسیر:} اگر نسبت درون‌نهشتی بالا باشد، تطبیق‌ها از نظر هندسی سازگارترند و احتمال درست بودن $\mathbf{H}$ بیشتر است.

\paragraph{نرخ موفقیت.}
برای آستانهٔ تحمل $\delta$ (پیکسل):
\[
\text{Success}(\delta)=\mathbb{I}(\text{MRE}<\delta),
\]
که در آن MRE می‌تواند «میانگین خطای بازفرافکنی» یا «میانه» باشد و $\mathbb{I}$ تابع نشانگر (۰/۱) است.

\subsection{شکل‌های لازم برای گزارش}
ارزیاب باید شکل‌های زیر را تولید کند:
\begin{itemize}
  \item \texttt{overlay\_grid.png}: بررسی کیفی هم‌پوشانی‌ها.
  \item \texttt{reproj\_error\_hist.png}: توزیع خطا (کشف دنباله‌های سنگین و شکست‌ها).
  \item \texttt{runtime\_breakdown.png}: گلوگاه‌های زمانی.
  \item \texttt{ablation\_plot.png}: حساسیت به پارامترها.
  \item \texttt{failure\_case\_*.png}: مرزهای اعتبار و تحلیل علت شکست.
\end{itemize}

\subsection{حداقل ابلیشن علمی}
برای ادعای پایداری، حداقل این ابلیشن‌ها پیشنهاد می‌شود:
\begin{itemize}
  \item \texttt{nfeatures} در چند مقدار،
  \item $\tau$ در چند مقدار،
  \item cross-check روشن/خاموش،
  \item CLAHE روشن/خاموش.
\end{itemize}
نتیجه‌گیری علمی باید بر «ناحیه‌های پایدار» تکیه کند، نه صرفاً یک نقطهٔ بهینه.
